\section{Conclusion}\label{sec:conclusion}

The central distinction is between
\[
\text{same endpoint judgement} \;\neq\; \text{same warrant}.
\]
Grounded transport carries evidence; composition and grounded-proper contexts preserve it; erasure yields ordinary substitutability and loses information irreversibly; located obstructions explain failed transport; open assessments preserve unresolved obligations, and a refutation does not discard them; and extraction retains the evidential data that downstream audits require.

Ordinary rewriting answers whether one expression may replace another under a relation. Grounded transport additionally answers through which process, under which lineage, within what coverage, and with which remaining obligations that replacement is warranted. Erasure recovers ordinary substitutability, but ordinary substitutability cannot reconstruct the warrant that was erased.

The calculus does not make declarations true. It makes precisely explicit what has been declared, checks what can be checked, and reports what has not been resolved.
